\chapter{Metrik Fonksiyonlar ve Sözleşmeler}

Bu bölümde metrik uzaylar arasındaki fonksiyon türleri (izometri, Lipschitz, sürekli)
ve bunların sınıflandırılması incelenmektedir.

%-------------------------------------------------------------------
\section{Konu}

\subsection{Metrik Fonksiyon Sınıfları}

$f: (M_1, d_1) \to (M_2, d_2)$ fonksiyonu:

\begin{table}[h]
\centering
\begin{tabular}{ll}
\toprule
\textbf{Sınıf} & \textbf{Koşul} \\
\midrule
Genişlemez (non-expansive) & $d_2(f(x),f(y)) \leq d_1(x,y)$ \\
Lipschitz & $\exists K>0:\; d_2(f(x),f(y)) \leq K \cdot d_1(x,y)$ \\
Üniform sürekli & $\forall\varepsilon>0\;\exists\delta>0:\; d_1(x,y)<\delta \Rightarrow d_2(f(x),f(y))<\varepsilon$ \\
Sürekli & Her noktada $\varepsilon$-$\delta$ sürekliliği \\
İzometri & $d_2(f(x),f(y)) = d_1(x,y)$ \\
Benzerlik & $\exists c>0:\; d_2(f(x),f(y)) = c \cdot d_1(x,y)$ \\
Homeomorfizma & Bijektif sürekli; ters de sürekli \\
\bottomrule
\end{tabular}
\caption{Metrik fonksiyon sınıfları}
\end{table}

\textbf{Sıralama:} İzometri $\Rightarrow$ Genişlemez $\Rightarrow$ Lipschitz
$\Rightarrow$ Üniform sürekli $\Rightarrow$ Sürekli.

\begin{figure}[ht]
\centering
\begin{tikzpicture}[font=\small]
  % --- domain (M1, d1) ---
  \draw[rounded corners=8pt, thick] (-0.4,-1.6) rectangle (3.2,1.8);
  \node at (-0.0,1.45) {$(M_1,d_1)$};
  \fill[blue!70!black] (0.6,0.9) circle (1.8pt) node[left, font=\scriptsize] {$x$};
  \fill[blue!70!black] (2.2,-0.9) circle (1.8pt) node[right, font=\scriptsize] {$y$};
  \draw[blue!70!black, thick] (0.6,0.9) -- node[above right, font=\scriptsize] {$d_1(x,y)$} (2.2,-0.9);
  % --- codomain (M2, d2) ---
  \draw[rounded corners=8pt, thick] (5.6,-1.6) rectangle (9.4,1.8);
  \node at (9.0,1.45) {$(M_2,d_2)$};
  \fill[violet!80!black] (6.6,0.9) circle (1.8pt) node[left, font=\scriptsize] {$f(x)$};
  \fill[violet!80!black] (8.2,-0.9) circle (1.8pt) node[right, font=\scriptsize] {$f(y)$};
  \draw[violet!80!black, thick] (6.6,0.9) -- node[above right, font=\scriptsize] {$d_2(f x,f y)$} (8.2,-0.9);
  % --- arrow f ---
  \draw[-{Stealth}, thick] (3.45,0.1) to[bend left=10] node[above, pos=0.5] {$f$} (5.35,0.1);
  % --- equality note ---
  \node[font=\scriptsize, align=center] at (4.4,-2.1)
    {izometri: $d_2(f x, f y) = d_1(x, y)$ (mesafe korunur)};
\end{tikzpicture}

\caption{İzometri: $f$, $x$ ile $y$ arasındaki mesafeyi korur ---
         $d_2(f(x),f(y)) = d_1(x,y)$.}
\end{figure}

\begin{sezgi}
Bir izometriyi, bir kâğıt parçasını \emph{buruşturmadan ve germeden} başka bir
yere taşıyan bir hareket gibi düşünün: döndürebilir, kaydırabilir,
yansıtabilirsiniz, ama iki nokta arasındaki mesafe asla değişmez. Düzlemde
öteleme, dönme ve yansıma izometridir. Benzerlik ise kâğıdı oranını koruyarak
büyüten/küçülten bir fotokopi makinesidir ($c=2 \Rightarrow$ her mesafe iki
katına çıkar); $c=1$ alındığında benzerlik tekrar izometriye iner.
\end{sezgi}

\subsection{Lipschitz Sabiti}

\[
K = \sup_{x \neq y} \frac{d_2(f(x),f(y))}{d_1(x,y)}.
\]

İzometri: $K=1$ ve eşitlik; benzerlik: $K=c$ sabit.

\begin{figure}[ht]
\centering
\begin{tikzpicture}[font=\small]
  % --- domain (M1, d1) ---
  \draw[rounded corners=8pt, thick] (-0.4,-1.7) rectangle (3.2,1.9);
  \node at (-0.0,1.55) {$(M_1,d_1)$};
  \fill[blue!70!black] (0.5,1.05) circle (1.8pt) node[left, font=\scriptsize] {$x$};
  \fill[blue!70!black] (2.4,-1.05) circle (1.8pt) node[right, font=\scriptsize] {$y$};
  \draw[blue!70!black, very thick] (0.5,1.05) -- node[above right, font=\scriptsize] {$d_1(x,y)$} (2.4,-1.05);
  % --- codomain (M2, d2): noktalar birbirine cok daha yakin ---
  \draw[rounded corners=8pt, thick] (5.6,-1.7) rectangle (9.4,1.9);
  \node at (9.0,1.55) {$(M_2,d_2)$};
  \fill[violet!80!black] (7.2,0.45) circle (1.8pt) node[above left, font=\scriptsize] {$f(x)$};
  \fill[violet!80!black] (7.9,-0.25) circle (1.8pt) node[below right, font=\scriptsize] {$f(y)$};
  \draw[violet!80!black, very thick] (7.2,0.45) -- node[above right, font=\scriptsize] {$\le k\,d_1(x,y)$} (7.9,-0.25);
  % --- arrow f ---
  \draw[-{Stealth}, thick] (3.45,0.05) to[bend left=10] node[above, pos=0.5] {$f$} (5.35,0.05);
  % --- contraction note ---
  \node[font=\scriptsize, align=center] at (4.4,-2.2)
    {büzülme: $d_2(f x, f y) \le k\, d_1(x,y),\; k<1$ (mesafe daralır)};
\end{tikzpicture}

\caption{Büzülme: $k<1$ olduğunda $f$ noktaları birbirine yaklaştırır ---
         $d_2(f(x),f(y)) \le k\,d_1(x,y)$.}
\end{figure}

\begin{sezgi}
Lipschitz sabiti $K$, fonksiyonun ``en dik eğimi'' gibidir: hiçbir nokta
çiftinde mesafeyi $K$ katından fazla büyütemez. $K<1$ olduğunda her adımda mesafe
kesin olarak küçülür --- buna \textbf{büzülme (contraction)} denir. Büzülmeyi
tekrar tekrar uyguladığınızda tüm noktalar tek bir noktaya (sabit noktaya) doğru
hunilenir; Banach sabit-nokta teoreminin kalbi budur.
\end{sezgi}

\begin{figure}[ht]
\centering
\begin{tikzpicture}[font=\small]
  % eksenler
  \draw[-{Stealth[length=2.2mm]}] (-0.3,0) -- (9.3,0) node[right] {$x$};
  \draw[-{Stealth[length=2.2mm]}] (0,-0.3) -- (0,3.4) node[above] {$f(x)$};
  % egim sinirli (Lipschitz) bir egri: tum noktalarda ayni delta is gorur
  \draw[blue!70!black, very thick] (0.2,0.4) .. controls (3.0,1.0) and (5.5,2.2) .. (9.0,2.9);
  % iki ornek nokta + ortak delta bandi
  \def\dx{0.9}
  % nokta A
  \draw[gray!55,dashed] (2.0,0) -- (2.0,1.05);
  \fill[red!75!black] (2.0,1.05) circle (1.6pt);
  \draw[orange!75!black,thick] (2.0-\dx,-0.18) -- (2.0+\dx,-0.18)
       node[midway,below,font=\scriptsize]{$\delta$};
  % nokta B (cok daha saglarda) -- AYNI delta
  \draw[gray!55,dashed] (7.0,0) -- (7.0,2.55);
  \fill[red!75!black] (7.0,2.55) circle (1.6pt);
  \draw[orange!75!black,thick] (7.0-\dx,-0.18) -- (7.0+\dx,-0.18)
       node[midway,below,font=\scriptsize]{$\delta$};
  % aciklama
  \node[font=\scriptsize, align=center, blue!70!black] at (5.0,3.25)
    {tüm noktalar için \emph{ortak} $\delta$ ($\varepsilon$ verildi)};
  \node[font=\scriptsize, align=center] at (4.6,-0.95)
    {düzgün süreklilik: $\delta$ noktadan bağımsız seçilebilir};
\end{tikzpicture}

\caption{Düzgün süreklilik: $\varepsilon$ verildiğinde aynı $\delta$ tüm
         noktalarda iş görür (noktadan bağımsız).}
\end{figure}

\begin{karsiornek}
Her sürekli fonksiyon düzgün (üniform) sürekli \emph{değildir}. $\mathbb{R}^+$
üzerinde $f(x) = 1/x$ süreklidir, ama düzgün sürekli değildir: $x$ sıfıra
yaklaştıkça eğim sınırsız diklenir, bu yüzden verilen bir $\varepsilon$ için
\emph{tek} bir $\delta$ bütün noktalarda iş göremez ($x \approx 0.001$ civarında
çok küçük bir $\delta$ gerekirken $x \approx 10$ civarında çok daha büyüğü
yeterlidir). Aynı şekilde $\mathbb{R}$ üzerinde $f(x)=x^2$ süreklidir ama düzgün
sürekli değildir. Düzgün süreklilik $\delta$'nın \emph{noktadan bağımsız}
seçilebilmesini ister; Lipschitz fonksiyonlar bunu $\delta=\varepsilon/K$ ile
otomatik sağlar.
\end{karsiornek}

%-------------------------------------------------------------------
\section{Teoremler}

\begin{theorem}
İzometri $\Rightarrow$ homeomorfizma.
\begin{proof}[İspat eskizi]
Önce \textbf{injektiflik}: $f$ bir izometri ve $f(x)=f(y)$ olsun. O zaman
$d_1(x,y) = d_2(f(x),f(y)) = 0$, dolayısıyla metriğin ayırma aksiyomundan $x=y$:
izometri daima injektiftir. Bijektif kabul edersek $f^{-1}$ vardır ve o da bir
izometridir. İzometri her mesafeyi koruduğu için $\varepsilon$-$\delta$
sürekliliğini $\delta=\varepsilon$ ile sağlar: hem $f$ hem $f^{-1}$ süreklidir.
Bijektif + sürekli + sürekli ters $=$ homeomorfizma.
\end{proof}
\end{theorem}

\begin{theorem}
Benzerlik, $c=1 \Rightarrow$ İzometri.
\end{theorem}

\begin{theorem}
Lipschitz $\Rightarrow$ Üniform sürekli $\Rightarrow$ Sürekli.
\begin{proof}[İspat eskizi]
$f$ sabiti $K>0$ olan Lipschitz olsun. Verilen $\varepsilon>0$ için
$\delta = \varepsilon/K$ seçin; $d_1(x,y)<\delta$ ise
$d_2(f(x),f(y)) \le K\,d_1(x,y) < K(\varepsilon/K) = \varepsilon$. Seçilen
$\delta$ noktaya bağlı değildir, dolayısıyla $f$ düzgün süreklidir; bu özel
olarak sürekliliği de içerir.
\end{proof}
\end{theorem}

\begin{theorem}
Kompakt metrik uzaydan metrik uzaya her sürekli fonksiyon üniform süreklidir.
\end{theorem}

%-------------------------------------------------------------------
\section{Algoritmalar}

\subsection{classify\_finite\_metric\_map — $O(|X|^2)$}

\begin{algorithm}
\caption{ClassifyMap$(X, d_1, Y, d_2, f)$}
\begin{algorithmic}[1]
\State $K \leftarrow 0$; isometry $\leftarrow$ True; similarity\_ratio $\leftarrow$ None
\For{each pair $(x_1, x_2)$ with $x_1 \neq x_2$}
    \State ratio $\leftarrow d_2(f(x_1),f(x_2)) / d_1(x_1,x_2)$
    \State $K \leftarrow \max(K, \text{ratio})$
    \If{$d_2(f(x_1),f(x_2)) \neq d_1(x_1,x_2)$}
        \State isometry $\leftarrow$ False
    \EndIf
    \If{similarity\_ratio is None}
        \State similarity\_ratio $\leftarrow$ ratio
    \ElsIf{ratio $\neq$ similarity\_ratio}
        \State similarity\_ratio $\leftarrow$ None
    \EndIf
\EndFor
\State \Return MetricMapProfile(lipschitz\_constant=$K$, isometry=isometry, $\ldots$)
\end{algorithmic}
\end{algorithm}

Karmaşıklık: $O(|X|^2)$.

%-------------------------------------------------------------------
\section{pytop API}

\begin{lstlisting}[language=Python]
from pytop.metric_spaces import FiniteMetricSpace
from pytop.metric_map_taxonomy import (
    classify_finite_metric_map,
    metric_map_profile,
    MetricMapProfile,
)
\end{lstlisting}

\texttt{classify\_finite\_metric\_map(fms1, fms2, f)} $\to$ \texttt{MetricMapProfile} döner (Result değil).

\texttt{MetricMapProfile} alanları: \texttt{isometry}, \texttt{similarity}, \texttt{similarity\_ratio},
\texttt{lipschitz\_constant}, \texttt{non\_expansive}, \texttt{bijective}, \texttt{lipschitz},
\texttt{continuous}, \texttt{uniformly\_continuous}, \texttt{homeomorphism}, \texttt{name}, \texttt{certification}.

%-------------------------------------------------------------------
\section{Örnekler}

\subsection*{Örnek 5.1 — Özdeşlik: İzometri}

\begin{lstlisting}[language=Python]
f_id = {1: 1, 2: 2, 3: 3, 4: 4}
prof_id = classify_finite_metric_map(fms, fms, f_id)
print("isometry:", prof_id.isometry)
print("lipschitz_constant:", prof_id.lipschitz_constant)
\end{lstlisting}

\begin{verbatim}
isometry: True
lipschitz_constant: 1.0
\end{verbatim}

\subsection*{Örnek 5.2 — 2× Ölçekleme: Benzerlik}

\begin{lstlisting}[language=Python]
f_scale = {1: 2, 2: 4, 3: 6, 4: 8}
prof_scale = classify_finite_metric_map(fms_line1, fms_line2, f_scale)
print("similarity:", prof_scale.similarity)
print("similarity_ratio:", prof_scale.similarity_ratio)
print("lipschitz_constant:", prof_scale.lipschitz_constant)
\end{lstlisting}

\begin{verbatim}
similarity: True
similarity_ratio: 2.0
lipschitz_constant: 2.0
\end{verbatim}

\subsection*{Örnek 5.3 — Sabit Fonksiyon: K=0}

\begin{lstlisting}[language=Python]
f_const = {1: 1, 2: 1, 3: 1, 4: 1}
prof_const = classify_finite_metric_map(fms, fms, f_const)
print("lipschitz_constant:", prof_const.lipschitz_constant)
print("bijective:", prof_const.bijective)
\end{lstlisting}

\begin{verbatim}
lipschitz_constant: 0.0
bijective: False
\end{verbatim}

\subsection*{Örnek 5.6 — Büzülme (Contraction): $K = 1/2$}

Öklid doğrusunda her mesafeyi yarıya indiren bir eşleme bir \textbf{büzülmedir}
($K<1$): Banach teoreminin sözleşme koşulunu sağlar.

\begin{lstlisting}[language=Python]
src = [0, 2, 4, 6]
img = [0, 1, 2, 3]
d_src = {(a, b): abs(a - b) for a in src for b in src}
d_img = {(a, b): abs(a - b) for a in img for b in img}
M_src = FiniteMetricSpace(carrier=tuple(src), distance=d_src)
M_img = FiniteMetricSpace(carrier=tuple(img), distance=d_img)
f_half = {0: 0, 2: 1, 4: 2, 6: 3}
prof_half = classify_finite_metric_map(M_src, M_img, f_half)
print("lipschitz_constant:", prof_half.lipschitz_constant)
print("similarity_ratio:", prof_half.similarity_ratio)
print("isometry:", prof_half.isometry)
\end{lstlisting}

\begin{verbatim}
lipschitz_constant: 0.5
similarity_ratio: 0.5
isometry: False
\end{verbatim}

\subsection*{Örnek 5.7 — Ayrık Metrikte Döngü: İzometri ve Homeomorfizma}

Ayrık metrikte her permütasyon mesafeyi korur: bir izometri ve (bijektif
olduğundan) homeomorfizmadır.

\begin{lstlisting}[language=Python]
dp = [1, 2, 3]
d_disc = {(a, b): (0 if a == b else 1) for a in dp for b in dp}
M_disc = FiniteMetricSpace(carrier=tuple(dp), distance=d_disc)
f_cycle = {1: 2, 2: 3, 3: 1}
prof_cycle = classify_finite_metric_map(M_disc, M_disc, f_cycle)
print("isometry:", prof_cycle.isometry)
print("homeomorphism:", prof_cycle.homeomorphism)
\end{lstlisting}

\begin{verbatim}
isometry: True
homeomorphism: True
\end{verbatim}

%-------------------------------------------------------------------
\section{Alıştırmalar}

\subsection*{Kodlama}

\begin{enumerate}[label=K\arabic*.]
\item $\{1,2,3\}$ üzerinde $f(1)=2, f(2)=3, f(3)=1$ (döngü) fonksiyonunu ayrık
      metrikte \texttt{classify\_finite\_metric\_map} ile sınıflandırın.
\item Öklid metrikli $\{0,1,2,3,4\}$ üzerinde $f(x) = x+1$ kaydırma fonksiyonunun
      Lipschitz sabitini hesaplayın.
\item Sabit fonksiyon $f(x) = 1$ için $K=0$ olduğunu doğrulayın.
\item Öklid doğrusunda $\{0,2,4,6\} \to \{0,1,2,3\}$ büzülmesini ($f(x)=x/2$) kurun;
      \texttt{lipschitz\_constant}, \texttt{similarity\_ratio}, \texttt{non\_expansive}
      alanlarını yazdırın ve $K<1$ olduğunu doğrulayın.
\item $\{1,2,3\}$ üzerinde özdeşlik olmayan bir döngüyü hem ayrık hem Öklid
      metrikte sınıflandırın; hangi metrikte izometri olduğunu \texttt{isometry}
      alanlarıyla karşılaştırın.
\end{enumerate}

\subsection*{Teori}

\begin{enumerate}[label=T\arabic*.]
\item İzometri $\Rightarrow$ genişlemez $\Rightarrow$ Lipschitz ($K=1$) zincirini ispatlayın.
\item Lipschitz $\Rightarrow$ üniform sürekli olduğunu $\delta = \varepsilon/K$ seçimi
      ile ispatlayın.
\item İzometrinin daima injektif olduğunu ispatlayın ($f(x)=f(y) \Rightarrow
      d_2(f(x),f(y))=0$); ardından düzlemde sabit bir fonksiyonun neden izometri
      olamadığını açıklayın.
\end{enumerate}
